%%
%% part2-book-components/ch12-tables.tex
%%
%% Chapter 12: Tables — Modern tables with tabularray,
%% the matintable wrapper, and RTL-aware column types.
%%

\chapter{جدول}
\label{chap:tables}

\section{چرا این موضوع مهم است؟}
\label{sec:tables-why}

جدول‌ها یکی از ابزارهای مهم در کتاب‌های فنی و
علمی هستند. از مقایسه‌ی الگوریتم‌ها تا نمایش
داده‌های تجربی، جدول‌ها به خواننده کمک می‌کنند
تا اطلاعات را سریع‌تر درک کند.

\lr{MatinBook} از بسته‌ی \lr{tabularray} برای
جدول‌ها استفاده می‌کند. این بسته، برخلاف
\lr{tabularx}، با \lr{xepersian} کاملاً سازگار
است و از حروف‌چینی راست‌به‌چپ پشتیبانی می‌کند.

\mnote{بسته‌ی \lr{tabularx} با \lr{xepersian} ناسازگار است و خطای \lr{Illegal pream-token} می‌دهد. هرگز از آن استفاده نکنید.}

\section{محیط \lr{matintable}}
\label{sec:tables-matintable}

ساده‌ترین راه برای ساخت جدول در \lr{MatinBook}،
استفاده از محیط \env{matintable} است. این
محیط، بخش‌های تکراری جدول را خودکار می‌کند:

\begin{latin}
\begin{minted}{latex}
\begin{matintable}{|\pc{عنوان جدول}|}{tab:my-table}
\begin{tblr}{
    width = \textwidth,
    colspec = {l l X[c] X[c]},
    hlines, vlines,
    colsep = 8pt,
    rowsep = 4pt,
    row{1} = {font=\bfseries},
}
|\pc{نوع}| & |\pc{فونت}| & |\pc{وضعیت}| & |\pc{توضیحات}| \\
|\pc{فارسی}| & \lr{XB Niloofar} & |\pc{فعال}| & |\pc{خوانایی بالا}| \\
|\pc{لاتین}| & \lr{Times New Roman} & |\pc{فعال}| & |\pc{استاندارد آکادمیک}| \\
\end{tblr}
\end{matintable}
\end{minted}
\end{latin}

خروجی:

\begin{matintable}{نمونه‌ی جدول با \lr{matintable}}{tab:sample}
\begin{tblr}{
    width = \textwidth,
    colspec = {l l X[c] X[c]},
    hlines, vlines,
    colsep = 8pt,
    rowsep = 4pt,
    row{1} = {font=\bfseries},
}
نوع & فونت & وضعیت & توضیحات \\
فارسی & \lr{XB Niloofar} & فعال & خوانایی بالا \\
لاتین & \lr{Times New Roman} & فعال & استاندارد آکادمیک \\
\end{tblr}
\end{matintable}

\mnote{دستور \lr{\textbackslash matintable} سه بخش تکراری را خودکار می‌کند: وسط‌چین کردن، کپشن، و برچسب.}

\section{ساختار \lr{tblr}}
\label{sec:tables-tblr}

محیط \env{tblr} داخل \env{matintable} قرار
می‌گیرد و تنظیمات جدول را تعیین می‌کند:

\begin{itemize}
    \item \lr{width = \textbackslash textwidth}:
    عرض کل جدول.
    \item \lr{colspec = \{...\}}: نوع ستون‌ها.
    \item \lr{hlines, vlines}: خطوط افقی و
    عمودی.
    \item \lr{colsep = 8pt}: فاصله‌ی داخلی
    ستون‌ها.
    \item \lr{rowsep = 4pt}: فاصله‌ی داخلی
    ردیف‌ها.
    \item \lr{row\{1\} = \{font=\textbackslash bfseries\}}:
    ردیف اول پررنگ.
\end{itemize}

\mnote{تنظیمات \lr{tblr} با کاما از هم جدا می‌شوند. ترتیب آن‌ها مهم نیست.}

\section{انواع ستون}
\label{sec:tables-columns}

\lr{MatinBook} از سه نوع ستون پشتیبانی می‌کند:

\begin{itemize}
    \item \lr{l}: ستون چپ‌چین (برای متن
    لاتین).
    \item \lr{c}: ستون وسط‌چین (برای اعداد
    یا نمادها).
    \item \lr{X[c]}: ستون خودکار با تراز
    وسط (برای متن فارسی).
\end{itemize}

\subsection{ستون \lr{X} در \lr{tabularray}}
\label{sec:tables-x-column}

ستون \lr{X} یک ستون \textbf{خودکار} است که
عرض آن به‌طور خودکار محاسبه می‌شود. می‌توانید
تراز آن را با \lr{[c]}، \lr{[l]} یا \lr{[r]}
تعیین کنید:

\begin{latin}
\begin{minted}{latex}
colspec = {l l X[c] X[c]}
\end{minted}
\end{latin}

\mnote{ستون \lr{X} فقط در \lr{tabularray} وجود دارد و در \lr{tabular} معمولی کار نمی‌کند.}

\subsection{ستون‌های \lr{p} (برای جدول‌های قدیمی)}
\label{sec:tables-p-columns}

\lr{MatinBook} سه نوع ستون \lr{p} هم ارائه می‌دهد
که برای \env{tabular} معمولی کار می‌کنند:

\begin{itemize}
    \item \lr{L\{width\}}: ستون چپ‌چین.
    \item \lr{R\{width\}}: ستون راست‌چین.
    \item \lr{C\{width\}}: ستون وسط‌چین.
\end{itemize}

\mnote{این ستون‌ها فقط برای سازگاری با کدهای قدیمی نگه داشته شده‌اند. برای جدول‌های جدید، از \lr{tblr} و ستون \lr{X} استفاده کنید.}

\section{عرض جدول}
\label{sec:tables-width}

برای اینکه جدول دقیقاً به عرض متن بچسبد، از
\lr{width = \textbackslash textwidth} استفاده
کنید:

\begin{latin}
\begin{minted}{latex}
\begin{tblr}{
    width = \textwidth,
    colspec = {l l X[c] X[c]},
    ...
}
\end{minted}
\end{latin}

\mnote{اگر از \lr{width} استفاده نکنید، عرض جدول به‌طور خودکار محاسبه می‌شود و ممکن است بیشتر از عرض متن شود.}

\subsection{جدول‌های باریک}
\label{sec:tables-narrow}

برای جدول‌های باریک (مثلاً ۲ ستون)، می‌توانید
از عرض کمتر استفاده کنید:

\begin{latin}
\begin{minted}{latex}
\begin{tblr}{
    width = 0.5\textwidth,
    colspec = {c c},
    ...
}
\end{minted}
\end{latin}

\section{خطوط جدول}
\label{sec:tables-rules}

\lr{tabularray} از سه نوع خط پشتیبانی می‌کند:

\begin{itemize}
    \item \lr{hlines}: خطوط افقی بین همه‌ی
    ردیف‌ها.
    \item \lr{vlines}: خطوط عمودی بین همه‌ی
    ستون‌ها.
    \item \lr{hlines, vlines}: هر دو.
\end{itemize}

\subsection{خطوط انتخابی}
\label{sec:tables-selective-rules}

می‌توانید خطوط را برای ردیف یا ستون خاصی تعیین
کنید:

\begin{latin}
\begin{minted}{latex}
\begin{tblr}{
    hline{1,Z} = {1pt},       % |\pc{خط ضخیم در بالا و پایین}|%
    hline{2} = {0.5pt},        % |\pc{خط نازک بعد از ردیف اول}|%
    ...
}
\end{minted}
\end{latin}

\mnote{\lr{Z} به معنای «آخرین ردیف» است. \lr{1,Z} یعنی ردیف اول و آخر.}

\section{رنگ در جدول}
\label{sec:tables-color}

\lr{tabularray} از رنگ‌آمیزی سلول‌ها پشتیبانی
می‌کند:

\begin{latin}
\begin{minted}{latex}
\begin{tblr}{
    row{1} = {bg=matinblue!20},    % |\pc{پس‌زمینه‌ی ردیف اول}|%
    cell{2}{3} = {bg=matingreen!20}, % |\pc{سلول خاص}|%
    ...
}
\end{minted}
\end{latin}

\mnote{برای رنگ‌آمیزی، از پالت \lr{MatinBook} استفاده کنید: \lr{matinblue}، \lr{matingreen}، \lr{matinorange} و...}

\subsection{مثال کامل با رنگ}
\label{sec:tables-color-example}

\begin{matintable}{مقایسه‌ی الگوریتم‌ها با رنگ}{tab:colored}
\begin{tblr}{
    width = \textwidth,
    colspec = {l X[c] X[c]},
    hlines, vlines,
    colsep = 8pt,
    rowsep = 4pt,
    row{1} = {font=\bfseries, bg=matinblue!20},
}
الگوریتم & پیچیدگی & کاربرد \\
مرتب‌سازی سریع & \lr{$O(n \log n)$} & \textcolor{matingreen}{سریع} \\
مرتب‌سازی ادغامی & \lr{$O(n \log n)$} & \textcolor{matingreen}{پایدار} \\
مرتب‌سازی حبابی & \lr{$O(n^2)$} & \textcolor{matinred}{کند} \\
\end{tblr}
\end{matintable}

\section{ادغام سلول‌ها}
\label{sec:tables-merge}

\lr{tabularray} از ادغام سلول‌ها پشتیبانی می‌کند:

\begin{latin}
\begin{minted}{latex}
\begin{tblr}{
    cell{1}{1} = {r=2}{c},   % |\pc{ادغام دو ردیف}|%
    cell{1}{2} = {c=2}{c},   % |\pc{ادغام دو ستون}|%
}
\end{minted}
\end{latin}

\mnote{در \lr{tabularray}، ادغام سلول‌ها با \lr{r} (ردیف) و \lr{c} (ستون) انجام می‌شود.}

\section{جدول‌های چندصفحه‌ای}
\label{sec:tables-multipage}

برای جدول‌های طولانی که بیش از یک صفحه هستند،
از \lr{longtblr} استفاده کنید:

\begin{latin}
\begin{minted}{latex}
\begin{longtblr}{
    width = \textwidth,
    colspec = {l l X[c]},
    hlines, vlines,
    rowhead = 1,
}
|\pc{نوع}| & |\pc{مقدار}| & |\pc{توضیحات}| \\
...
\end{longtblr}
\end{minted}
\end{latin}

\mnote{گزینه‌ی \lr{rowhead = 1} باعث می‌شود ردیف اول (سرصفحه) در هر صفحه‌ی جدید تکرار شود.}

\section{جدول‌های فارسی و لاتین}
\label{sec:tables-persian-latin}

در جدول‌های \lr{MatinBook}، می‌توانید فارسی و
لاتین را با هم ترکیب کنید:

\begin{latin}
\begin{minted}{latex}
\begin{matintable}{|\pc{ترکیب فارسی و لاتین}|}{tab:mixed}
\begin{tblr}{
    width = \textwidth,
    colspec = {l X[c]},
    hlines, vlines,
    colsep = 8pt,
    rowsep = 4pt,
    row{1} = {font=\bfseries},
}
|\pc{مفهوم}| & |\pc{معادل لاتین}| \\
|\pc{الگوریتم}| & \lr{Algorithm} \\
|\pc{برنامه‌نویسی}| & \lr{Programming} \\
\end{tblr}
\end{matintable}
\end{minted}
\end{latin}

خروجی:

\begin{matintable}{ترکیب فارسی و لاتین}{tab:mixed}
\begin{tblr}{
    width = \textwidth,
    colspec = {l X[c]},
    hlines, vlines,
    colsep = 8pt,
    rowsep = 4pt,
    row{1} = {font=\bfseries},
}
مفهوم & معادل لاتین \\
الگوریتم & \lr{Algorithm} \\
برنامه‌نویسی & \lr{Programming} \\
\end{tblr}
\end{matintable}

\mnote{در جدول‌ها، متن لاتین باید داخل \lr{\textbackslash lr\{...\}} قرار گیرد. متن فارسی نیازی به \lr{\textbackslash lr} ندارد.}

\section{جدول در حالت تمام‌عرض}
\label{sec:tables-fullwidth}

اگر جدول شما پهن است و به فضای بیشتری نیاز
دارد، می‌توانید از عرض کامل متن استفاده کنید:

\begin{latin}
\begin{minted}{latex}
\begin{tblr}{
    width = \textwidth,
    colspec = {X[l] X[c] X[c] X[c] X[c]},
    ...
}
\end{minted}
\end{latin}

\mnote{عرض \lr{\textbackslash textwidth} در متن اصلی حدود \lr{۱۰.۳} سانتی‌متر است. در پیش‌متن و پس‌متن، این عرض بیشتر است.}

\section{ارجاع به جدول}
\label{sec:tables-references}

برای ارجاع به جدول، از \cmd{cref} استفاده کنید:

\begin{latin}
\begin{minted}{latex}
|\pc{جدول}|\cref{tab:sample}|\pc{نمونه‌ای از جدول‌های}|\lr{MatinBook}|\pc{است.}|%
\end{minted}
\end{latin}

خروجی: جدول \cref{tab:sample} نمونه‌ای از جدول‌های
\lr{MatinBook} است.

\mnote{نام‌گذاری برچسب جدول‌ها با پیشوند \lr{tab:} انجام می‌شود: \lr{tab:sample}.}

\section{خطاهای رایج}
\label{sec:tables-mistakes}

\subsection{استفاده از \lr{tabularx}}
\label{sec:tables-mistake-tabularx}

\lr{tabularx} با \lr{xepersian} ناسازگار است و
خطای زیر را می‌دهد:

\begin{latin}
\begin{minted}{text}
Package array: Illegal pream-token
  (\TX@col@width): `c' used.
\end{minted}
\end{latin}

\textbf{راه‌حل:} به‌جای \lr{tabularx} از
\lr{tabularray} استفاده کنید.

\subsection{عرض ثابت در ستون}
\label{sec:tables-mistake-fixed-width}

اگر از عرض ثابت (مثل \lr{C\{3cm\}}) استفاده
کنید، ممکن است جدول سرریز کند.

\textbf{راه‌حل:} از ستون \lr{X[c]} استفاده کنید
که عرض آن خودکار است.

\subsection{فراموش کردن \cmd{lr} در سلول‌های لاتین}
\label{sec:tables-mistake-lr}

اگر متن لاتین در سلول جدول را بدون \cmd{lr}
بنویسید، به‌صورت نامفهوم نمایش داده می‌شود:

\begin{latin}
\begin{minted}{latex}
% |\pc{اشتباه}|:
|\pc{الگوریتم}| & Algorithm \\

% |\pc{درست}|:
|\pc{الگوریتم}| & \lr{Algorithm} \\
\end{minted}
\end{latin}

\section{تمرین‌ها}
\label{sec:tables-exercises}

\begin{exercise}
    یک جدول دو ستونی بسازید که عرض آن دقیقاً
    به عرض متن بچسبد.
    \label{exc:tables-basic}
\end{exercise}

\begin{exercise}
    یک جدول با چهار ستون بسازید که دو ستون
    آن لاتین و دو ستون فارسی باشد.
    \label{exc:tables-mixed}
\end{exercise}

\begin{exercise}
    یک جدول با رنگ‌آمیزی ردیف اول و یک سلول
    خاص بسازید.
    \label{exc:tables-color}
\end{exercise}

\begin{exercise}
    یک جدول طولانی بسازید که به دو صفحه
    تقسیم شود و از \lr{longtblr} استفاده کنید.
    \label{exc:tables-long}
\end{exercise}

\section{خلاصه}
\label{sec:tables-summary}

در این فصل، با جدول در \lr{MatinBook} آشنا
شدیم:

\begin{itemize}
    \item \lr{MatinBook} از \lr{tabularray}
    استفاده می‌کند که با \lr{xepersian} سازگار
    است.

    \item محیط \env{matintable} بخش‌های تکراری
    جدول را خودکار می‌کند.

    \item محیط \env{tblr} تنظیمات جدول را تعیین
    می‌کند: عرض، ستون‌ها، خطوط، فاصله‌ها.

    \item سه نوع ستون: \lr{l}، \lr{c}، \lr{X[c]}.

    \item برای عرض دقیق، از
    \lr{width = \textbackslash textwidth} استفاده
    کنید.

    \item خطوط جدول با \lr{hlines} و \lr{vlines}
    تعیین می‌شوند.

    \item رنگ‌آمیزی با \lr{bg=...} انجام می‌شود.

    \item جدول‌های طولانی با \lr{longtblr} ساخته
    می‌شوند.

    \item خطاهای رایج شامل استفاده از
    \lr{tabularx}، عرض ثابت، و فراموش کردن
    \cmd{lr} است.
\end{itemize}

در فصل بعدی (\cref{chap:margin-notes})، با
یادداشت حاشیه‌ای در \lr{MatinBook} آشنا
می‌شوید و یاد می‌گیرید که چگونه از فضای
حاشیه‌ی پهن استفاده کنید.